\chapter{Классическая информация, обратимая динамика и матрицы}

\lectureoverview{Квантовая механика требует от нас рассуждать за пределами области, в которой формировалась обычная интуиция. Поэтому мы начнем на шаг раньше — с классической информации: битов, конфигураций и обратимых законов обновления. То же обсуждение естественным образом приведет нас к векторам и матрицам, которые дадут нам язык для квантовых битов.}

\section{Интуиция, сформированная не для того мира}

Наша физическая интуиция стара. Задолго до того, как кто-либо записал уравнение, животное, преследующее добычу, должно было оценивать относительную скорость и направление. Если лев догоняет антилопу, он реагирует на изменяющееся расстояние между ними и, следовательно, неявным образом — на относительную скорость. Нервные системы, хорошо справлявшиеся с такими оценками, имели преимущество; те, что справлялись с ними плохо, оставляли мало потомков. Значительная часть того, что мы называем здравым смыслом в вопросах движения, сформировалась в этом обычном мире.

Это наследие полезно, но ограниченно. В обычной жизни ничто не движется со скоростью, близкой к скорости света, поэтому эволюция не снабдила нас инстинктом для релятивистского пространства-времени. Физик сначала изучает абстрактную геометрию и лишь позже вырабатывает для нее новую интуицию. С квантовой механикой еще труднее. Привычное нам представление о частице происходит из опыта с камнями и стрелами: в каждый момент времени объект имеет положение, последовательность положений образует траекторию, а изменения положения определяют скорость. Электрон вовсе не обязан укладываться в эту картину. Неопределенность положения и отсутствие классической траектории не имеют аналога в опыте, сформировавшем наше восприятие.

В том, что интуиция здесь подводит, нет никакого парадокса. Парадокс возник бы лишь в том случае, если бы инстинкты, приспособленные к камням, животным и умеренным скоростям, каким-то образом заранее знали об электронах. Правильный ответ — построить иную интуицию на основе небольшого набора точных правил. Здесь акцент будет сделан не на обычном пути через волновую механику и уравнение Schr\"odinger. Мы сосредоточимся на логической структуре квантовой механики и на квантовой информации. Прежде чем спрашивать, что такое квантовый бит, мы должны в точности понять, как работает обычный бит.

\section{Биты и количество информации}

Физика неотделима от информации. Иногда эта информация выражается действительным числом, но простейший возможный элемент информации — это ответ на вопрос, у которого есть только две альтернативы: да или нет, орел или решка, вверх или вниз.

\begin{definition}
Классический \emph{бит} — это физический вопрос с ровно двумя возможными ответами.
\end{definition}

Мы можем обозначить эти ответы как \(0\) и \(1\). Сами обозначения не несут внутреннего смысла; подошли бы любые два различных символа. Поскольку обозначения Дирака вскоре станут незаменимыми, введем их уже сейчас, хотя пока они избыточны:
\begin{equation}
  \ket{0},
  \qquad
  \ket{1}.
\end{equation}
Вертикальная черта и угловая скобка образуют \emph{кет}, правую половину бра--кета. Ничего квантового не появилось лишь потому, что мы сменили обозначения.

Для нескольких различимых битов запишем ответы в фиксированном порядке. Например, шесть различимых монет могут иметь конфигурацию
\begin{equation}
  \ket{0\,1\,1\,0\,1\,0}.
\end{equation}
Запятые необязательны, когда порядок уже понятен. У одного бита 2 возможных состояния, у двух битов — 4, и каждый добавленный бит удваивает это число. Таким образом, система из \(n\) битов имеет
\begin{equation}
  N_s=2^n
  \label{eq:l01-state-count}
\end{equation}
состояний. Обратно,
\begin{equation}
  n=\log_2 N_s.
  \label{eq:l01-bits-from-states}
\end{equation}

Уравнение~\eqref{eq:l01-bits-from-states} подсказывает определение, полезное даже тогда, когда число состояний не является степенью двойки:
\begin{equation}
  I=\log_2 N_s.
  \label{eq:l01-information}
\end{equation}
Следовательно, шестигранная кость несет \(\log_2 6\) бит информации. Результат не обязан быть целым числом. Системы, буквально собранные из двоичных переключателей, будут удобными примерами, но логарифм — более общая мера.

\section{Кодирование кажущегося непрерывным мира}

Биты не ограничиваются монетами. Любую измеряемую величину, известную с конечной точностью, можно записать как конечную двоичную строку. Предположим, например, что температуру в комнате требуется знать лишь с выбранным числом значащих цифр. Выразите измеренное значение в системе счисления с основанием 2 и сохраните цифры, необходимые для этой точности. Большая точность требует более длинной строки, но принцип при этом не меняется.

Аналогично можно поступить с полем, меняющимся в пространстве. Разбейте комнату на мелкую упорядоченную решетку ячеек, запишите значение поля в каждой ячейке с конечной точностью и соедините полученные двоичные строки в одну. Уменьшение размеров ячеек улучшает пространственное разрешение; сохранение большего числа двоичных разрядов улучшает разрешение по значению поля. Температурные поля, электрические поля и магнитные поля можно таким способом закодировать с любой заданной конечной точностью.

Конфигурации частиц дают еще более буквальную двоичную картину. Разбейте пространство на достаточно малые ячейки, чтобы каждая содержала не более одной частицы. Зададим
\begin{equation}
  n_i=
  \begin{cases}
    1,&\text{cell \(i\) is occupied},\\
    0,&\text{cell \(i\) is empty}.
  \end{cases}
  \label{eq:l01-occupancy-bit}
\end{equation}
Полная конфигурация — это упорядоченная строка
\begin{equation}
  \ket{n_1n_2n_3\cdots}.
\end{equation}
Это очень длинный ответ на последовательность вопросов типа да или нет: есть ли частица в первой ячейке, во второй и так далее? Это физическая версия игры \emph{Двадцать вопросов}, только число вопросов может быть огромным.

Эта конструкция объясняет, почему цифровые компьютеры могут моделировать физику. Компьютер хранит конечные строки, а конечное разрешение превращает измеряемые величины и поля именно в такие строки. Приближение можно систематически улучшать, уточняя решетку, шаг времени или число сохраняемых цифр.

\begin{classroomqa}
\lecturetimestamp{00:29:43}
\audiencequestion{Является ли точность представления температуры произвольной?}
\lecturerresponse{Она выбирается в соответствии с тем приближением, которое нам нужно. Конечная строка дает конечную точность; сохранение большего числа двоичных разрядов дает более точное приближение. Здесь не утверждается, что конечная строка точно задает произвольное классическое действительное число.}
\end{classroomqa}

\begin{classroomqa}
\lecturetimestamp{00:30:11}
\audiencequestion{Что происходит для общего действительного числа, особенно иррационального?}
\lecturerresponse{Точное действительное число может потребовать бесконечно много битов. Некоторые двоичные разложения рациональных чисел периодичны, а разложение типичного иррационального числа никогда не обрывается. Реальный компьютер всегда хранит лишь конечное число цифр. Суть в том, что можно закодировать любую заданную конечную точность, а не в том, что каждое классическое действительное число имеет конечное точное описание.}
\end{classroomqa}

\section{Состояние и закон движения}

До сих пор мы описывали систему в один момент времени. Для движения нужна вторая составляющая.

\begin{definition}
\emph{Конфигурация} — это полное состояние системы в заданный момент времени. Закон движения в дискретном времени — это правило, которое отображает каждую конфигурацию в конфигурацию на один шаг времени позже.
\end{definition}

Это различие элементарно, но существенно. Снимок распределения электронов — это конфигурация; это еще не орбита. Если пространство заменить решеткой, а время — последовательностью \(t_0,t_1,t_2,\ldots\), то движение есть последовательность
\begin{equation}
  \ket{s(t_0)}
  \longmapsto
  \ket{s(t_1)}
  \longmapsto
  \ket{s(t_2)}
  \longmapsto\cdots.
  \label{eq:l01-discrete-motion}
\end{equation}
Поэтому физической модели нужны и множество возможных состояний, и правило обновления.

\begin{classroomqa}
\lecturetimestamp{00:31:12}
\audiencequestion{Как это описание через ячейки описывает электрон, движущийся по орбите?}
\lecturerresponse{Сначала задайте только конфигурацию в некоторый момент: какие ячейки заняты, при фиксированном ядре, если это входит в модель. Затем дискретизируйте время. Орбитальное движение — это правило, которое преобразует одну строку занятости в следующую. Правило может быть сложным из-за упорядочения ячеек, но это все равно обновление информации.}
\end{classroomqa}

Полезно изобразить абстрактное \emph{пространство состояний}. Каждая точка в этом пространстве обозначает целую конфигурацию, а не положение в обычном пространстве. Решетка занятости \(10\times10\) содержит 100 битов и, следовательно, \(2^{100}\) возможных конфигураций.\editorialnote{В устной лекции при разборе этого примера кратко говорится \(10^{100}\), но число битов фиксирует подразумеваемое значение как \(2^{100}\).} Каждая из этих конфигураций соответствует одной точке в пространстве состояний.

Для одного бита есть только две точки. Тождественный закон оставляет обе неизменными:
\begin{equation}
  0\mapsto0,
  \qquad
  1\mapsto1.
  \label{eq:l01-identity-update}
\end{equation}
Другая обратимая возможность меняет их местами:
\begin{equation}
  0\mapsto1,
  \qquad
  1\mapsto0.
  \label{eq:l01-flip-update}
\end{equation}
Это флип-флоп. Физический пример, эволюционирующий сам по себе, не столь очевиден, как выключатель света, переключаемый вручную, потому что в этом случае пришлось бы расширить систему. День и ночь дают приближенный двухсостояний цикл; маятник тоже чередует крайние положения, хотя его полное состояние включает и все промежуточные положения, и скорости.

\begin{classroomqa}
\lecturetimestamp{00:40:38}
\audiencequestion{Мог бы обычный выключатель света быть примером закона флип-флоп?}
\lecturerresponse{Нет, если его переключает человек, потому что тогда выключатель не эволюционирует сам по себе; человек входит в расширенную динамическую систему. Нам нужна замкнутая система, собственный закон которой попеременно меняет ее состояние. День и ночь подсказывают эту идею, а маятник напоминает, что у реального примера обычно больше двух состояний.}
\end{classroomqa}

Двухбитная система имеет четыре конфигурации. Один закон обновления может проводить все четыре через один общий цикл. Другой может состоять из двух отдельных двухциклов. Это петли в логическом пространстве конфигураций, а не обязательно петли в физическом пространстве. Стрелки задают закон: из какого бы состояния ни начиналась система, следование по стрелкам определяет её последующую историю.

\section{Обратимость и сохранение различий}

Не всякая диаграмма со стрелками имеет один и тот же физический статус. Рассмотрим правило, при котором два различных состояния сливаются в одно и то же последующее состояние:
\begin{equation}
  a\mapsto c,
  \qquad
  b\mapsto c.
  \label{eq:l01-merging-update}
\end{equation}
Следующее состояние определено и при старте из \(a\), и при старте из \(b\), но после достижения \(c\) мы не можем определить, какое прошлое имело место. Различие между \(a\) и \(b\) исчезло.

\begin{figure}[htbp]
\centering
\includegraphics[width=0.82\linewidth]{lecture_01_figure_02.png}
\caption{Во время обсуждения обратимости на доске остаются несколько диаграмм обновления. Замкнутые циклы соседствуют со схемой слияния, а рядом с ними записан квантовый термин \emph{унитарность}.}
\lectureframe{00:53:21}
\end{figure}

\begin{figure}[htbp]
\centering
\begin{tikzpicture}[>=Latex,thick,node distance=1.15cm]
  \node[circle,fill,inner sep=1.8pt,label=above:\(1\)] (r1) {};
  \node[circle,fill,inner sep=1.8pt,right=of r1,label=above:\(2\)] (r2) {};
  \node[circle,fill,inner sep=1.8pt,below=of r2,label=below:\(3\)] (r3) {};
  \node[circle,fill,inner sep=1.8pt,left=of r3,label=below:\(4\)] (r4) {};
  \draw[->] (r1) -- (r2);
  \draw[->] (r2) -- (r3);
  \draw[->] (r3) -- (r4);
  \draw[->] (r4) -- (r1);
  \node[below=0.5cm of r4,anchor=north west] {\small reversible cycle};

  \node[circle,fill,inner sep=1.8pt,right=4.2cm of r1,label=above:\(a\)] (a) {};
  \node[circle,fill,inner sep=1.8pt,below=1.25cm of a,label=below:\(b\)] (b) {};
  \node[circle,fill,inner sep=1.8pt,right=1.6cm of a,yshift=-0.625cm,label=above:\(c\)] (c) {};
  \node[circle,fill,inner sep=1.8pt,right=of c,label=above:\(d\)] (d) {};
  \draw[->] (a) -- (c);
  \draw[->] (b) -- (c);
  \draw[->] (c) -- (d);
  \node[below=0.55cm of b,anchor=north west] {\small information-losing merge};
\end{tikzpicture}
\caption{Обратимое правило сопоставляет каждому состоянию одного предшественника и одного преемника. Правило со слиянием имеет единственное будущее, но неоднозначное прошлое.}
\end{figure}

Фундаментальная классическая динамика не допускает такой потери, если сохраняются все степени свободы. Требуемое свойство имеет две стороны:
\begin{equation}
  \boxed{\begin{gathered}
    \text{one unique future and one unique past}\\
    \text{for every state}
  \end{gathered}.}
  \label{eq:l01-reversibility}
\end{equation}
Следовательно, для конечного пространства состояний обновление должно быть перестановкой. У каждого состояния есть ровно одна исходящая стрелка и ровно одна входящая, так что диаграмма распадается на непересекающиеся циклы.

\begin{classroomqa}
\lecturetimestamp{00:46:32}
\audiencequestion{Чем отличается диаграмма со слиянием?}
\lecturerresponse{Она не состоит целиком из петель. Более существенно, если система оказалась в слитом состоянии, это не говорит нам, какой у неё был предшественник. Будущее мы можем построить, а прошлое восстановить не можем. Один бит различимости прошлого утрачен.}
\end{classroomqa}

Это понятие — \emph{обратимость во времени}, а не обязательно точная симметрия относительно обращения времени. Цикл по часовой стрелке имеет определённую ориентацию вперёд. Исходный закон не переводит систему по нему против часовой стрелки, однако существует обратный закон, который это делает. Точная симметрия относительно обращения времени задаёт дополнительный вопрос: получаются ли те же самые динамические законы после применения соответствующей операции обращения? Сама по себе обратимость спрашивает лишь о том, можно ли однозначно воспроизвести историю в обратном направлении.

\begin{classroomqa}
\lecturetimestamp{00:49:58}
\audiencequestion{Разве единственность прошлого — это просто симметрия относительно обращения времени?}
\lecturerresponse{Не совсем. Стрелки цикла могут быть направлены только в одну сторону, так что у закона есть ориентация. Обратимость означает, что, обратив каждую стрелку, можно построить однозначно определённый обратный закон. Само по себе это не означает, что обратный закон тождествен исходному.}
\end{classroomqa}

И не следует чинить необратимую диаграмму, помещая в слитом состоянии две исходящие стрелки. Это сделало бы обратное движение неоднозначным по-другому: системе понадобилось бы вероятностное предписание, например выбирать одну ветвь в \(30\%\) случаев, а другую — в \(70\%\). Мы по-прежнему строим детерминистическую классическую динамику, в которой одно настоящее состояние определяет одно будущее и одно прошлое.

Всякая реальная система в конечном счёте квантовомеханична. Почему большие объекты кажутся классическими и почему их квантовые альтернативы становится трудно наблюдать, — вопрос более поздний. Соответствующий квантовый принцип сохранит то же логическое ядро. В полушутку можно было бы назвать однозначность прошлого и будущего \emph{уникарностью}. Её настоящее название — \emph{унитарность}: квантовую эволюцию можно однозначно прослеживать вперёд или назад.

\section{Кажущаяся потеря информации и огрубление описания}

Макроскопический опыт, по-видимому, противоречит обратимости. Поместите все молекулы газа в левый угол герметично закрытой комнаты и отпустите их. Вскоре газ заполнит комнату. Если же начать с правого угла, итоговая макроскопическая картина будет выглядеть так же. Слились ли эти два различных начальных состояния?

Лишь после того, как мы отбросили детали. Микроскопическое состояние задаёт каждую молекулу по отдельности. Если бы можно было с достаточной точностью проследить каждое положение и каждую скорость, обращение этих движений вспять позволило бы различить две подготовки. На практике такое вычисление непомерно трудно, а требуемая точность быстро растёт, но практическая невозможность не равносильна принципиальному стиранию. Макроскопическое утверждение \emph{воздух заполняет комнату} отождествляет огромное число различных микроскопических конфигураций с одним огрублённым состоянием.

Именно здесь вступает в дело второй закон термодинамики. Мы теряем способность различать детальные состояния, потому что не отслеживаем все микроскопические переменные. Кажется, что информация исчезает, потому что наше описание огрублено, а не потому, что детерминистический закон в точном описании сливает различные точки.

\begin{classroomqa}
\lecturetimestamp{00:56:11}
\audiencequestion{Потребовало бы отслеживание всех молекул в классической комнате бесконечно многих битов?}
\lecturerresponse{Для точного классического описания — да. Положения и скорости являются действительными числами, и прослеживание хаотической многочастичной системы на всё больших временах требует всё большей точности. Квантовая механика меняет этот вывод: для замкнутой области с ограниченной энергией физически различимое описание может быть дискретным, а не произвольным списком точных классических действительных чисел.}
\end{classroomqa}

Вопрос не только академический. Может ли чёрная дыра уничтожать информацию, десятилетиями было предметом спора; Stephen Hawking поначалу утверждал, что да. В этом споре ставится вопрос, является ли кажущееся слияние всего лишь ещё одним огрублённым описанием или подлинным нарушением квантовой унитарности. Мы не будем решать его, пока квантовая механика не предоставит точный язык, но он показывает, как много опирается на простые диаграммы со стрелками.

\section{Строки, столбцы и внутренние произведения}

Прежде чем заменять классические биты квантовыми, нам понадобится немного линейной алгебры. Для наших целей вектор можно рассматривать просто как конечный упорядоченный список действительных чисел. Вектор с четырьмя компонентами можно записать горизонтально,
\begin{equation}
  A^{\mathsf T}=(A_1,A_2,A_3,A_4),
  \label{eq:l01-row-vector}
\end{equation}
или вертикально,
\begin{equation}
  A=
  \begin{pmatrix}
    A_1\\ A_2\\ A_3\\ A_4
  \end{pmatrix}.
  \label{eq:l01-column-vector}
\end{equation}
Слово \emph{четырёхмерный} означает лишь то, что есть четыре компоненты. Нет необходимости представлять стрелку в обычном пространстве. Горизонтальный объект — это вектор-строка, а вертикальный — вектор-столбец.

\begin{figure}[htbp]
\centering
\includegraphics[width=0.92\linewidth]{lecture_01_figure_03.png}
\caption{Одни и те же компоненты записаны как строка и как столбец. На этом этапе это действительные числа, так что комплексного сопряжения нет.}
\lectureframe{01:03:05}
\end{figure}

\begin{classroomqa}
\lecturetimestamp{01:03:14}
\audiencequestion{Разве вектор-строка не является эрмитовым сопряжением вектора-столбца?}
\lecturerresponse{Это будет правильным языком, как только компоненты станут комплексными. Тогда транспонирование должно сопровождаться комплексным сопряжением. Пока элементы действительны, нам нужно лишь отличать запись в строку от записи в столбец.}
\end{classroomqa}

Пусть
\begin{equation}
  B=(B_1,B_2,B_3,B_4)
\end{equation}
— строка с тем же числом элементов, что и у столбца \(A\). Их произведение определяется попарным сопоставлением позиций и суммированием:
\begin{align}
  BA
  &=
  (B_1,B_2,B_3,B_4)
  \begin{pmatrix}
    A_1\\ A_2\\ A_3\\ A_4
  \end{pmatrix}\\
  &=B_1A_1+B_2A_2+B_3A_3+B_4A_4.
  \label{eq:l01-inner-product}
\end{align}
Результат — число, а не ещё один вектор и не матрица. В абстрактном векторном пространстве эта операция называется внутренним произведением.

\begin{classroomqa}
\lecturetimestamp{01:06:44}
\audiencequestion{Это скалярное произведение?}
\lecturerresponse{Да. Для обычных векторов в трёхмерном пространстве это привычное скалярное произведение. Название \emph{внутреннее произведение} подчёркивает, что одна и та же операция умножения строки на столбец применяется в векторных пространствах любой конечной размерности.}
\end{classroomqa}

\section{Матрицы как машины для обновления состояний}

Матрица — это квадратный массив чисел, который действует на вектор и порождает другой вектор. Думайте о ней как о машине: подайте на вход старый столбец, примените фиксированное правило и получите новый столбец. В четырёх измерениях,
\begin{equation}
  M=
  \begin{pmatrix}
    M_{11}&M_{12}&M_{13}&M_{14}\\
    M_{21}&M_{22}&M_{23}&M_{24}\\
    M_{31}&M_{32}&M_{33}&M_{34}\\
    M_{41}&M_{42}&M_{43}&M_{44}
  \end{pmatrix}.
  \label{eq:l01-general-matrix}
\end{equation}
Первый индекс помечает строку, а второй — столбец. На \(M\) можно смотреть либо как на набор столбцов, либо как на набор строк; обе точки зрения будут полезны.

Чтобы вычислить \(A'=MA\), возьмите внутреннее произведение каждой строки \(M\) со входным столбцом:
\begin{equation}
  A'_i=(MA)_i=\sum_{j=1}^{4}M_{ij}A_j.
  \label{eq:l01-matrix-vector}
\end{equation}
Например,
\begin{align}
  A'_1&=M_{11}A_1+M_{12}A_2+M_{13}A_3+M_{14}A_4,\\
  A'_2&=M_{21}A_1+M_{22}A_2+M_{23}A_3+M_{24}A_4.
\end{align}
Третья и четвёртая компоненты получаются по тому же образцу.

\begin{figure}[htbp]
\centering
\includegraphics[width=0.94\linewidth]{lecture_01_figure_04.png}
\caption{Матрица действует на вектор-столбец построчно. Первые две выходные компоненты явно расписаны на доске.}
\lectureframe{01:12:45}
\end{figure}

\subsection{Кодирование классической конфигурации}

Предположим, что система имеет четыре возможных состояния. Представим их столбцами с одной единицей
\begin{equation}
  e_1=
  \begin{pmatrix}1\\0\\0\\0\end{pmatrix},
  \quad
  e_2=
  \begin{pmatrix}0\\1\\0\\0\end{pmatrix},
  \quad
  e_3=
  \begin{pmatrix}0\\0\\1\\0\end{pmatrix},
  \quad
  e_4=
  \begin{pmatrix}0\\0\\0\\1\end{pmatrix}.
  \label{eq:l01-one-hot}
\end{equation}
Положение единственной \(1\) указывает, какая классическая конфигурация реализована.

Рассмотрим циклическое правило обновления
\begin{equation}
  1\longrightarrow2\longrightarrow3\longrightarrow4\longrightarrow1.
  \label{eq:l01-four-cycle}
\end{equation}
Оно представляется матрицей
\begin{equation}
  P=
  \begin{pmatrix}
    0&0&0&1\\
    1&0&0&0\\
    0&1&0&0\\
    0&0&1&0
  \end{pmatrix}.
  \label{eq:l01-permutation-matrix}
\end{equation}
В самом деле,
\begin{equation}
  \begin{aligned}
    Pe_1&=e_2, & Pe_2&=e_3,\\
    Pe_3&=e_4, & Pe_4&=e_1.
  \end{aligned}
\end{equation}
Построение на доске исправляется по ходу записи: сначала с пятью состояниями, затем с четырьмя. В ~\eqref{eq:l01-permutation-matrix} зафиксировано окончательное циклическое правило, а не отброшенный промежуточный вариант.

Та же матрица действует на произвольный столбец:
\begin{equation}
  P
  \begin{pmatrix}A\\B\\C\\D\end{pmatrix}
  =
  \begin{pmatrix}D\\A\\B\\C\end{pmatrix}.
  \label{eq:l01-cyclic-shift}
\end{equation}
Для классического состояния только одна компонента равна \(1\), но если разрешить произвольные элементы, становится видна лежащая в основе линейная операция. В классических детерминистических обновлениях используются специальные нуль-единичные матрицы, в каждой строке и каждом столбце которых стоит ровно одна \(1\). В квантовой механике матричная эволюция сохранится, но множество допустимых элементов сильно расширится.

\subsection{Композиция — это умножение матриц}

Если один шаг по времени задается соотношением
\begin{equation}
  V'=MV,
\end{equation}
то второе применение дает
\begin{equation}
  V''=MV'=M(MV)=M^2V.
  \label{eq:l01-two-updates}
\end{equation}
Следовательно, умножение матриц — не произвольное алгебраическое упражнение. Оно задает композицию физических обновлений. Пять шагов по времени представляются как \(M^5\).

Для двух матриц \(2\times2\)
\begin{equation}
  M=
  \begin{pmatrix}
    M_{11}&M_{12}\\
    M_{21}&M_{22}
  \end{pmatrix},
  \qquad
  N=
  \begin{pmatrix}
    N_{11}&N_{12}\\
    N_{21}&N_{22}
  \end{pmatrix},
\end{equation}
произведение \(C=MN\) имеет элементы
\begin{equation}
  C_{ij}=\sum_k M_{ik}N_{kj}.
  \label{eq:l01-matrix-product}
\end{equation}
В частности,
\begin{align}
  C_{11}&=M_{11}N_{11}+M_{12}N_{21},\\
  C_{12}&=M_{11}N_{12}+M_{12}N_{22},\\
  C_{21}&=M_{21}N_{11}+M_{22}N_{21},\\
  C_{22}&=M_{21}N_{12}+M_{22}N_{22}.
\end{align}
Каждый элемент получается как произведение строки левой матрицы на столбец правой. Эквивалентно, можно подействовать \(M\) на каждый столбец \(N\) и поставить полученные столбцы рядом.

\begin{figure}[htbp]
\centering
\includegraphics[width=0.86\linewidth]{lecture_01_figure_05.png}
\caption{Умножение матриц, организованное как последовательное умножение матрицы--на столбец. Первая строка и столбцы второй матрицы отмечены, чтобы показать, откуда берется каждый элемент произведения.}
\lectureframe{01:26:30}
\end{figure}

Вектор-строка тоже может действовать слева. Если \(B=(B_1,\ldots,B_4)\), то \(BM\) — еще один вектор-строка, причем
\begin{equation}
  (BM)_j=\sum_i B_iM_{ij}.
  \label{eq:l01-row-times-matrix}
\end{equation}
И здесь правило то же: строка, умноженная на столбец. Форма записи показывает, что получается:
\begin{equation}
\begin{array}{c@{\qquad}c}
\text{row}\times\text{column}&=\text{number},\\
\text{matrix}\times\text{column}&=\text{column},\\
\text{row}\times\text{matrix}&=\text{row},\\
\text{matrix}\times\text{matrix}&=\text{matrix}.
\end{array}
\label{eq:l01-shape-rules}
\end{equation}

Эти операции полезно отработать на небольших численных примерах \(2\times2\) и \(3\times3\). Наряду с небольшим количеством математического анализа умножение строки--на столбец — основной новый математический навык, необходимый для предстоящей квантовой механики.

\begin{classroomqa}
\lecturetimestamp{01:31:50}
\audiencequestion{Какое ограничение обратимость накладывает на матрицу обновления?}
\lecturerresponse{У нее должна существовать обратная матрица. Обратная матрица реализует обратное обновление и возвращает каждое состояние к своему единственному предшественнику. Не всякая матрица обратима, так что алгебра воспроизводит различие, проведенное ранее с помощью диаграмм со стрелками. Для циклического примера \(P^{-1}=P^3\) и \(P^4=\mathbb{1}\).}
\end{classroomqa}

\section{Зачем разбирать алгебру до кубита?}

В первый час было дано качественное описание детерминистической классической физики в терминах битов и обратимых стрелок. Дальнейшая математика переводит это описание в форму, которая сохранится при переходе к квантовой механике. Конфигурация становится столбцом; закон на один шаг становится матрицей; многократная эволюция во времени становится произведением матриц; обратимость становится существованием обратной матрицы.

Мы еще не использовали квантовое правило. Эта задержка преднамеренна. Далее мы заменим классические альтернативы \(\ket{0}\) и \(\ket{1}\) кубитом и спросим, какие новые состояния возможны. Векторы-строки, векторы-столбцы, скалярные произведения и матрицы, введенные здесь, затем покажут, почему квантовая информация столь резко отличается от классической битовой строки.

\begin{classroomqa}
\lecturetimestamp{01:34:00}
\audiencequestion{Зачем вводить всю эту матричную алгебру, прежде чем сказать, для чего она нужна?}
\lecturerresponse{Сначала стоял качественный вопрос: как следует представлять детерминистическую физику с помощью состояний и законов обновления? Эта алгебра — подготовка к следующему шагу. Когда будут введены кубиты, эти операции уже не будут казаться отвлеченными; они станут естественным языком квантовых состояний и их эволюции.}
\end{classroomqa}